\chapter{Karta RV32I}
\label{app:karta}

Kodowanie wszystkich 37 instrukcji obliczeniowych, pamięciowych i sterujących RV32I
(bez \code{fence}, \code{ecall}, \code{ebreak}). Tabela jest wygenerowana z tablic
asemblera kursu (\plik{common/tools/rvtool.py}), który koduje instrukcje bit w bit
tak samo jak LLVM. Formaty i immediate'y: rys.~\ref{fig:formaty}.

% WYGENEROWANE z common/tools/rvtool.py — nie edytuj ręcznie
\begin{center}
\footnotesize
\renewcommand{\arraystretch}{1.0}
\begin{tabular}{@{}l c l l l l@{}}
\toprule
\textbf{Instr.} & \textbf{Fmt} & \textbf{opcode} & \textbf{funct3} & \textbf{funct7} & \textbf{Działanie} \\
\midrule
\code{lui} & U & \code{0110111} & \code{---} & \code{---} & rd = imm $\ll$ 12 \\
\code{auipc} & U & \code{0010111} & \code{---} & \code{---} & rd = pc + (imm $\ll$ 12) \\
\code{jal} & J & \code{1101111} & \code{---} & \code{---} & rd = pc+4; pc += imm \\
\code{jalr} & I & \code{1100111} & \code{000} & \code{---} & rd = pc+4; pc = (rs1+imm) \& \textasciitilde 1 \\
\midrule
\code{beq} & B & \code{1100011} & \code{000} & \code{---} & rs1 == rs2 \\
\code{bne} & B & \code{1100011} & \code{001} & \code{---} & rs1 != rs2 \\
\code{blt} & B & \code{1100011} & \code{100} & \code{---} & rs1 $<_s$ rs2 \\
\code{bge} & B & \code{1100011} & \code{101} & \code{---} & rs1 $\geq_s$ rs2 \\
\code{bltu} & B & \code{1100011} & \code{110} & \code{---} & rs1 $<_u$ rs2 \\
\code{bgeu} & B & \code{1100011} & \code{111} & \code{---} & rs1 $\geq_u$ rs2 \\
\midrule
\code{lb} & I & \code{0000011} & \code{000} & \code{---} & rd = sext(M[rs1+imm][7:0]) \\
\code{lh} & I & \code{0000011} & \code{001} & \code{---} & rd = sext(M[rs1+imm][15:0]) \\
\code{lw} & I & \code{0000011} & \code{010} & \code{---} & rd = M[rs1+imm] \\
\code{lbu} & I & \code{0000011} & \code{100} & \code{---} & rd = zext(M[rs1+imm][7:0]) \\
\code{lhu} & I & \code{0000011} & \code{101} & \code{---} & rd = zext(M[rs1+imm][15:0]) \\
\code{sb} & S & \code{0100011} & \code{000} & \code{---} & M[rs1+imm][7:0] = rs2[7:0] \\
\code{sh} & S & \code{0100011} & \code{001} & \code{---} & M[rs1+imm][15:0] = rs2[15:0] \\
\code{sw} & S & \code{0100011} & \code{010} & \code{---} & M[rs1+imm] = rs2 \\
\midrule
\code{addi} & I & \code{0010011} & \code{000} & \code{---} & rd = rs1 + imm \\
\code{slti} & I & \code{0010011} & \code{010} & \code{---} & rd = rs1 $<_s$ imm \\
\code{sltiu} & I & \code{0010011} & \code{011} & \code{---} & rd = rs1 $<_u$ imm \\
\code{xori} & I & \code{0010011} & \code{100} & \code{---} & rd = rs1 \^{} imm \\
\code{ori} & I & \code{0010011} & \code{110} & \code{---} & rd = rs1 | imm \\
\code{andi} & I & \code{0010011} & \code{111} & \code{---} & rd = rs1 \& imm \\
\code{slli} & I* & \code{0010011} & \code{001} & \code{0000000} & rd = rs1 $\ll$ shamt \\
\code{srli} & I* & \code{0010011} & \code{101} & \code{0000000} & rd = rs1 $\gg_u$ shamt \\
\code{srai} & I* & \code{0010011} & \code{101} & \code{0100000} & rd = rs1 $\gg_s$ shamt \\
\midrule
\code{add} & R & \code{0110011} & \code{000} & \code{0000000} & rd = rs1 + rs2 \\
\code{sub} & R & \code{0110011} & \code{000} & \code{0100000} & rd = rs1 - rs2 \\
\code{sll} & R & \code{0110011} & \code{001} & \code{0000000} & rd = rs1 $\ll$ rs2[4:0] \\
\code{slt} & R & \code{0110011} & \code{010} & \code{0000000} & rd = rs1 $<_s$ rs2 \\
\code{sltu} & R & \code{0110011} & \code{011} & \code{0000000} & rd = rs1 $<_u$ rs2 \\
\code{xor} & R & \code{0110011} & \code{100} & \code{0000000} & rd = rs1 \^{} rs2 \\
\code{srl} & R & \code{0110011} & \code{101} & \code{0000000} & rd = rs1 $\gg_u$ rs2[4:0] \\
\code{sra} & R & \code{0110011} & \code{101} & \code{0100000} & rd = rs1 $\gg_s$ rs2[4:0] \\
\code{or} & R & \code{0110011} & \code{110} & \code{0000000} & rd = rs1 | rs2 \\
\code{and} & R & \code{0110011} & \code{111} & \code{0000000} & rd = rs1 \& rs2 \\
\bottomrule
\end{tabular}
\end{center}


{\small I* --- format I, w którym \code{imm[11:5]} pełni rolę \code{funct7}, a
\code{imm[4:0]} to \code{shamt}. Wszystkie immediate'y poza U są rozszerzane
znakiem. \code{sext}/\code{zext}: rozszerzenie znakiem/zerami.}

\section*{Mapa pamięci maszyny kursu}

\begin{center}
\small
\begin{tabular}{@{}lll@{}}
\toprule
\code{0x0000\_0000}--\code{0x0000\_FFFF} & RAM 64 KiB & program, dane, stos \\
\code{0x1000\_0000} & TOHOST (W) & 1 = PASS, $(n \ll 1)|1$ = FAIL $n$ \\
\code{0x1000\_0004} & PUTCHAR (W) & znak \code{wdata[7:0]} \\
\code{0x1000\_0008} & CYCLES (R) & licznik cykli \\
\bottomrule
\end{tabular}
\end{center}

\chapter{Ściąga SystemVerilog}
\label{app:sv}

\section*{Składnia}
\begin{sv}
module nazwa import rv32i_pkg::*; #(parameter int W = 8) (
  input  logic         clk, rst,
  input  logic [W-1:0] a,
  output logic [W-1:0] y
);
  logic [W-1:0] t;                      // sygnał wewnętrzny
  localparam logic [3:0] STALA = 4'hA;  // stała lokalna

  assign t = a ^ {W{1'b1}};             // logika kombinacyjna (krótka)

  always_comb begin                     // logika kombinacyjna (dłuższa)
    y = '0;                             // wartość domyślna: brak zatrzasku
    case (t[1:0])
      2'd0:    y = a;
      default: y = t;
    endcase
  end

  always_ff @(posedge clk)              // rejestr
    if (rst) q <= '0;
    else     q <= y;
endmodule
\end{sv}

\section*{Literały i operatory}
\begin{center}
\small
\begin{tabularx}{\linewidth}{@{}l X@{}}
\toprule
\code{8'hFF}, \code{4'b1010}, \code{32'd7}, \code{'0}, \code{'1} & literały z rozmiarem; same zera/jedynki \\
\code{x[7:0]}, \code{x[8*i +: 8]} & wycinek stały; wycinek o zmiennym początku \\
\code{\{a, b\}}, \code{\{4\{b\}\}} & sklejanie; powielanie \\
\code{\& | \^{} \~{}} & bitowe; jako unarne: redukcja (\code{|x} = „jakikolwiek bit”) \\
\code{\&\& || !} & logiczne (1 bit) \\
\code{== !=}, \code{=== !==} & równość; dosłowna (z X i Z), tylko w testbenchach \\
\code{<< >> >>>} & przesunięcia; \code{>>>} arytmetyczne tylko dla signed \\
\code{\$signed(x)}, \code{\$unsigned(x)} & interpretacja ze znakiem / bez znaku \\
\code{W'(x)} & rzutowanie na szerokość \code{W} \\
\code{\$clog2(N)} & $\lceil \log_2 N \rceil$ (szerokość adresu dla $N$ elementów) \\
\bottomrule
\end{tabularx}
\end{center}

\section*{Zasady kursu}
\begin{enumerate}
  \item \code{always\_comb} + \code{=}, \code{always\_ff} + \code{<=}. Nigdy odwrotnie.
  \item W \code{always\_comb} każdy sygnał ma wartość domyślną na początku bloku.
  \item Każdy \code{case} ma \code{default}.
  \item Jeden zegar, reset synchroniczny aktywny w stanie 1.
  \item Stałe z pakietu, nie „magiczne liczby”.
  \item Uważaj na szerokości: \code{make lint} musi być czysty.
\end{enumerate}

\section*{Testbench}
\begin{sv}
`timescale 1ns/1ps
module tb_x;
  `include "tb_common.svh"
  logic clk = 0;
  always #5 clk = ~clk;
  initial begin
    tb_start();
    @(negedge clk);                       // zmieniaj wejścia na zboczu opadającym
    `CHECK_EQ(y, 8'h42, "opis")
    repeat (10) @(negedge clk);
    $display("a=%h b=%0d", a, b);
    tb_finish();
  end
endmodule
\end{sv}

\chapter{Ściąga poleceń}
\label{app:make}

\begin{center}
\small
\begin{tabularx}{\linewidth}{@{}l X@{}}
\toprule
\textbf{Polecenie} & \textbf{Działanie} \\
\midrule
\code{make test} & wszystkie testy modułu \\
\code{make unit [T=nazwa]} & testy jednostkowe (\plik{tb/tb\_nazwa.sv}) \\
\code{make wave T=nazwa} & test + przebiegi w GTKWave \\
\code{make lint} & Verilator \code{--lint-only -Wall} \\
\code{make synth T=moduł} & Yosys: komórki i najdłuższa ścieżka \\
\code{make progs [PAD=n] [SIM=verilator]} & 12 programów testowych na rdzeniu \\
\code{make prog P=08\_fib} & jeden program: wyjście, ślad \plik{build/P.trace}, listing \\
\code{make wave P=08\_fib} & przebiegi wykonania programu \\
\code{make cosim} & porównanie śladu z emulatorem \\
\code{make fuzz N=100 [LEN=300 SEED=1]} & losowe programy + porównanie \\
\code{make emu-test [P=...]} & (moduł 04) programy na emulatorze \\
\code{make run P=hello [CORE=08]} & (moduł 07) program w C na rdzeniu \\
\code{make clean} & usuwa \plik{build/} \\
\bottomrule
\end{tabularx}
\end{center}

\section*{Narzędzia w \plik{common/tools}}
\begin{sh}
python3 rvtool.py asm prog.S -o prog.hex -l prog.lst [-I katalog] [--pad-nops N]
python3 rvtool.py dis prog.hex          # albo: dis 0x00a50533
python3 rvtool.py bin2hex prog.bin -o prog.hex
python3 tracecmp.py ref.trace dut.trace
python3 rvgen.py --seed 7 -n 300 -o rand.S [--mext]
\end{sh}

\chapter{Słowniczek}
\label{app:slownik}

\begin{center}
\small
\begin{tabularx}{\linewidth}{@{}l l X@{}}
\toprule
\textbf{Polski} & \textbf{Angielski} & \textbf{Znaczenie} \\
\midrule
przerzutnik & flip-flop (FF) & element pamiętający bit na zboczu zegara \\
zatrzask & latch & element pamiętający sterowany poziomem; w RTL zwykle błąd \\
ścieżka krytyczna & critical path & najdłuższa droga kombinacyjna między rejestrami \\
bank rejestrów & register file & 32 rejestry z portami odczytu i zapisu \\
ścieżka danych & datapath & bloki przetwarzające dane (ALU, rejestry, multipleksery) \\
jednostka sterująca & control unit & dekoder wytwarzający sygnały sterujące \\
potok & pipeline & podział wykonania na etapy oddzielone rejestrami \\
hazard & hazard & sytuacja, w której potok dałby zły wynik bez dodatkowej logiki \\
przekazywanie & forwarding / bypass & dostarczenie wyniku prosto z rejestru potoku \\
wstrzymanie & stall & zatrzymanie części potoku na takt \\
bańka & bubble & pusty etap potoku (\code{valid = 0}) \\
unieważnienie & flush & zamiana instrukcji z błędnej ścieżki w bańki \\
zatwierdzenie & commit / retire & zakończenie instrukcji, widoczne w śladzie \\
stała natychmiastowa & immediate & stała zakodowana w instrukcji \\
rozszerzenie znakiem & sign extension & kopiowanie bitu znaku na starsze bity \\
co-symulacja & co-simulation & porównywanie RTL z modelem w trakcie wykonania \\
model referencyjny & golden model & niezależna implementacja traktowana jako wzorzec \\
pokrycie & coverage & miara tego, które zachowania zostały sprawdzone \\
\bottomrule
\end{tabularx}
\end{center}
